\documentclass[9pt]{article}

% Packages
\usepackage[margin=0.5in]{geometry}
\usepackage{amsmath, amssymb, amsthm}
\usepackage{enumitem}
\usepackage[most]{tcolorbox}
\usepackage{hyperref}
\usepackage{mathtools}
\usepackage{multicol}

% Dark axiom box style
\newtcolorbox{axiombox}{
  colback=black!85!white,
  coltext=white,
  colframe=black!90!white,
  boxrule=0pt,
  arc=4pt,
  left=8pt, right=8pt, top=6pt, bottom=6pt,
}

\begin{document}

% Title
\begin{center}
  {\LARGE \textbf{Gradient Computations}} \\[6pt]
  {\normalsize EECS 127} \\[10pt]
\end{center}

\vspace{4pt}

% Problem Statement
\begin{axiombox}
\textbf{Problem.} \; Compute the gradient $\nabla g(\mathbf{x})$ for each of the following functions, where $A \in \mathbb{R}^{m \times n}$, $\mathbf{b} \in \mathbb{R}^m$, and $\mathbf{x} \in \mathbb{R}^n$.

\begin{enumerate}[label=(\alph*), itemsep=3pt]
  \item $g_1(\mathbf{x}) = \mathbf{x}^T A \mathbf{x}$, where $A \in \mathbb{R}^{n \times n}$.
  \item $g_6(\mathbf{x}) = e^{\|\mathbf{x}\|_2^2}$.
  \item $g_7(\mathbf{x}) = e^{\|A\mathbf{x} - \mathbf{b}\|_2^2}$.
\end{enumerate}
\end{axiombox}

\vspace{8pt}

\begin{multicols}{2}

%% ---- Part (a) ---- %%
\section*{Part (a): $\nabla g_1(\mathbf{x})$}

We want $\nabla_{\mathbf{x}}(\mathbf{x}^T A \mathbf{x})$.

Expand the scalar $\mathbf{x}^T A \mathbf{x}$ in index notation:
\[
  g_1(\mathbf{x}) = \sum_{i,j} x_i \, A_{ij} \, x_j.
\]

Differentiating with respect to $x_k$:
\[
  \frac{\partial g_1}{\partial x_k}
  = \sum_j A_{kj}\, x_j + \sum_i x_i\, A_{ik}.
\]

In vector form, the first sum gives the $k$-th entry of $A\mathbf{x}$ and the second gives the $k$-th entry of $A^T \mathbf{x}$. Therefore:
\[
  \boxed{\nabla g_1(\mathbf{x}) = (A + A^T)\,\mathbf{x}.}
\]

\textbf{Note:} If $A$ is symmetric ($A = A^T$), this simplifies to $2A\mathbf{x}$.

\columnbreak

%% ---- Part (b) ---- %%
\section*{Part (b): $\nabla g_6(\mathbf{x})$}

Write $g_6(\mathbf{x}) = e^{h(\mathbf{x})}$ where $h(\mathbf{x}) = \|\mathbf{x}\|_2^2 = \mathbf{x}^T \mathbf{x}$.

By the chain rule for scalar compositions:
\[
  \nabla g_6 = e^{h(\mathbf{x})} \cdot \nabla h(\mathbf{x}).
\]

Since $\nabla(\mathbf{x}^T \mathbf{x}) = 2\mathbf{x}$:
\[
  \boxed{\nabla g_6(\mathbf{x}) = 2\,e^{\|\mathbf{x}\|_2^2}\,\mathbf{x}.}
\]

\vspace{8pt}

%% ---- Part (c) ---- %%
\section*{Part (c): $\nabla g_7(\mathbf{x})$}

Write $g_7(\mathbf{x}) = e^{h(\mathbf{x})}$ where $h(\mathbf{x}) = \|A\mathbf{x} - \mathbf{b}\|_2^2$.

First compute $\nabla h$. Let $\mathbf{r} = A\mathbf{x} - \mathbf{b}$, so $h = \mathbf{r}^T \mathbf{r}$. By the chain rule:
\[
  \nabla_{\mathbf{x}} h
  = \frac{\partial \mathbf{r}}{\partial \mathbf{x}}^T \nabla_{\mathbf{r}}(\mathbf{r}^T\mathbf{r})
  = A^T (2\mathbf{r})
  = 2A^T(A\mathbf{x} - \mathbf{b}).
\]

Applying the outer chain rule:
\[
  \boxed{\nabla g_7(\mathbf{x}) = 2\,e^{\|A\mathbf{x}-\mathbf{b}\|_2^2}\,A^T(A\mathbf{x}-\mathbf{b}).}
\]

\end{multicols}

\end{document}
